\section*{Contents}
\begin{tabularx}{0.98\linewidth}{@{}>{\bfseries}p{0.28\linewidth}X@{}}
  1 & Open Weights Go Architectural --- 公開の重みは構造で競う \\
  2 & Agent Coding Plane --- 道具の面が組み替わる \\
  3 & Flagship Receipts --- 看板の主張と受け止め方 \\
  4 & Infra \& Regional --- 動かす経済と地域の公開 \\
  5 & Safety \& Governance --- 安全の証拠と統治の対照 \\
  6 & Week in Review --- 変化・意味・次の観測点 \\
\end{tabularx}

\vspace{4mm}
\section*{This Week in AI}

\begin{tcolorbox}[enhanced,colback=SurveySoft,colframe=SurveyRule,boxrule=0.4pt,arc=1mm,left=3mm,right=3mm,top=2mm,bottom=2mm]
\textbf{公開の重みが構造で並んだ。}\quad 大きなMoEの公開と予告が同じ週に現れたが、動作規模の小ささも利用条件もモデルごとに確かめる必要がある。ひと括りの効率物語にはしない。
\end{tcolorbox}

\begin{tcolorbox}[enhanced,colback=SurveySoft,colframe=SurveyRule,boxrule=0.4pt,arc=1mm,left=3mm,right=3mm,top=2mm,bottom=2mm]
\textbf{道具の面が組み替わった。}\quad 続きの持ち運び、置き場の取り込み、会話の文脈化、審査と見張りが並び、便利さと預け先の問いが同時に前に出た。
\end{tcolorbox}

\begin{tcolorbox}[enhanced,colback=SurveySoft,colframe=SurveyRule,boxrule=0.4pt,arc=1mm,left=3mm,right=3mm,top=2mm,bottom=2mm]
\textbf{看板と土台と統治を分けて読む。}\quad 旗艦の測定値は出どころと留保ごと受け止め、動かす土台は出荷記録で固め、安全の証拠と二つの統治の形は対照として並べる。
\end{tcolorbox}

\vspace{2mm}
\begin{claimboundary}[この号の読み方]
通常本文の基準時点は2026-08-28 18:00 America/New\_York。本号で参照する公開資料が示す範囲を超えて断定しない。提供元や報道が伝える数値はそれぞれの主体に帰属させ、条件の違う測定値を横断的なランキングへ言い換えない。重み置き場や一次発表での確認を欠く仕様や利用条件は事実として扱わない。
\end{claimboundary}

\vspace{2mm}
\noindent\textit{読み解く軸：} 公開の構造効率と道具面の組み替えを、旗艦の看板の留保、動かす土台の出荷記録、安全の証拠と統治の対照とともに確かめることである。
